\begin{proof}[Proof of \cref{obj:lem:private-kernel-comparison}]
\begin{enumerate}
\item For adjacent datasets \(D,D'\) and a measurable output event \(E\), \cref{obj:def:private-kernels,obj:ass:pure-privacy} give
\[
M(D,E)\le e^\epsilon M(D',E),
\qquad
M(D',E)\le e^\epsilon M(D,E).
\]
Since the event probabilities belong to \([0,1]\),
\[
\begin{aligned}
M(D,E)-M(D',E)
&\le (e^\epsilon-1)M(D',E)\le e^\epsilon-1,\\
M(D',E)-M(D,E)
&\le (e^\epsilon-1)M(D,E)\le e^\epsilon-1.
\end{aligned}
\]
Convexity of the exponential, together with \(0<\epsilon\le1\) and \(e<3\), yields
\[
e^\epsilon
\le (1-\epsilon)e^0+\epsilon e^1,
\qquad
e^\epsilon-1\le (e-1)\epsilon\le2\epsilon.
\]
Consequently,
\[
|M(D,E)-M(D',E)|\le2\epsilon
\qquad\text{if }d_{\mathrm H}(D,D')=1.
\]
% lean: CausalSmith.Stat.PrivateCateRoughdesign.exp_privacy_increment_le
% lean: CausalSmith.Stat.PrivateCateRoughdesign.private_adjacent_event_bound

\item Fix a measurable output event \(E\), and define
\[
\varphi_E:\mathcal O^n\longrightarrow[0,1],
\qquad
\varphi_E(D)=M(D,E).
\]
We prove by induction on a finite set of coordinate positions that
\[
\mathcal J\subseteq\{1,\ldots,n\},
\quad
D_i=D'_i\ \text{for all }i\notin\mathcal J
\quad\Longrightarrow\quad
|\varphi_E(D)-\varphi_E(D')|
\le2\epsilon|\mathcal J|.
\]
For \(\mathcal J=\varnothing\), the datasets coincide and both sides of the inequality are zero.

For the induction step, suppose \(D,D'\) agree outside
\(\mathcal J\cup\{\iota\}\), where \(\iota\notin\mathcal J\). Define the intermediate dataset by
\[
(D_{\mathrm{mid}})_{\iota'}
=
\begin{cases}
D'_\iota,&\iota'=\iota,\\
D_{\iota'},&\iota'\ne\iota,
\end{cases}
\qquad \iota'\in\{1,\ldots,n\}.
\]
Then \(D_{\mathrm{mid}}\) and \(D'\) agree outside \(\mathcal J\), so the induction hypothesis gives
\[
|\varphi_E(D_{\mathrm{mid}})-\varphi_E(D')|
\le2\epsilon|\mathcal J|.
\]
If \(D_\iota=D'_\iota\), then \(D_{\mathrm{mid}}=D\) and their event-probability difference is zero. Otherwise, their Hamming distance is one, and the adjacent bound applies. In either case,
\[
|\varphi_E(D)-\varphi_E(D_{\mathrm{mid}})|\le2\epsilon.
\]
The triangle inequality therefore gives
\[
\begin{aligned}
|\varphi_E(D)-\varphi_E(D')|
&\le
|\varphi_E(D)-\varphi_E(D_{\mathrm{mid}})|
+
|\varphi_E(D_{\mathrm{mid}})-\varphi_E(D')|\\
&\le2\epsilon+2\epsilon|\mathcal J|
=2\epsilon|\mathcal J\cup\{\iota\}|.
\end{aligned}
\]
Apply this conclusion to
\[
\mathcal J=\{i\in\{1,\ldots,n\}:D_i\ne D'_i\}.
\]
It follows that, for every pair of datasets,
\[
|M(D,E)-M(D',E)|
\le2\epsilon d_{\mathrm H}(D,D').
\]
The empty-set case also covers \(n=0\).
% lean: CausalSmith.Stat.PrivateCateRoughdesign.event_bound_of_hamming_path
% lean: CausalSmith.Stat.PrivateCateRoughdesign.private_hamming_event_bound

\item Let \(\gamma\) be any coupling of \(Q,Q'\); explicitly,
\[
\gamma\in\operatorname{Prob}(\mathcal O^n\times\mathcal O^n),
\qquad
\gamma(F\times\mathcal O^n)=Q(F),
\qquad
\gamma(\mathcal O^n\times F)=Q'(F)
\]
for every measurable dataset event \(F\). Since Hamming distance is bounded by the dataset size,
\[
0\le
\int d_{\mathrm H}(D,D')\,\gamma(dD,dD')
\le
\int n\,\gamma(dD,dD')
=n.
\]
Thus the coupling cost is finite. The measurable function \(\varphi_E\) is bounded by one, so its integrals against the marginals and against \(\gamma\) are well defined. The marginal identities give
\[
(MQ)(E)-(MQ')(E)
=
\int\bigl(\varphi_E(D)-\varphi_E(D')\bigr)\,
\gamma(dD,dD').
\]
Using the integral triangle inequality and the preceding event bound,
\[
\begin{aligned}
|(MQ)(E)-(MQ')(E)|
&\le
\int|\varphi_E(D)-\varphi_E(D')|\,\gamma(dD,dD')\\
&\le
\int2\epsilon d_{\mathrm H}(D,D')\,\gamma(dD,dD')\\
&=
2\epsilon\int d_{\mathrm H}(D,D')\,\gamma(dD,dD').
\end{aligned}
\]
Taking the supremum over measurable output events proves
\[
d_{\mathrm{TV}}(MQ,MQ')
\le
2\epsilon\int d_{\mathrm H}(D,D')\,\gamma(dD,dD').
\]
% lean: CausalSmith.Stat.PrivateCateRoughdesign.private_TV_le_coupling_cost
This holds for every coupling. Since \(2\epsilon\) is positive and finite, taking the infimum gives
\[
d_{\mathrm{TV}}(MQ,MQ')
\le
\inf_{\gamma\ \mathrm{couples}\ Q,Q'}
2\epsilon\int d_{\mathrm H}(D,D')\,d\gamma
=
2\epsilon W_{\mathrm H}(Q,Q').
\]
% lean: CausalSmith.Stat.PrivateCateRoughdesign.private_kernel_comparison

\item Let
\[
K:\mathcal O^n\rightsquigarrow\mathcal B
\]
be an everywhere probability-valued Markov kernel. For a measurable output event \(E\), define
\[
\varphi_{K,E}:\mathcal O^n\longrightarrow[0,1],
\qquad
\varphi_{K,E}(D)=K(D,E).
\]
This function is measurable and bounded. Its layer representation under each of the probability laws is
\[
\begin{aligned}
\int\varphi_{K,E}(D)\,Q(dD)
&=\int_{(0,1]}Q\{D:\eta_{\mathrm{layer}}\le\varphi_{K,E}(D)\}\,d\eta_{\mathrm{layer}},\\
\int\varphi_{K,E}(D)\,Q'(dD)
&=\int_{(0,1]}Q'\{D:\eta_{\mathrm{layer}}\le\varphi_{K,E}(D)\}\,d\eta_{\mathrm{layer}}.
\end{aligned}
\]
These identities follow by integrating
\(\varphi_{K,E}(D)=\int_{(0,1]}\mathbf 1\{\eta_{\mathrm{layer}}\le\varphi_{K,E}(D)\}\,d\eta_{\mathrm{layer}}\).
Every displayed level set is measurable, and hence
\[
\left|
Q\{D:\eta_{\mathrm{layer}}\le\varphi_{K,E}(D)\}
-
Q'\{D:\eta_{\mathrm{layer}}\le\varphi_{K,E}(D)\}
\right|
\le d_{\mathrm{TV}}(Q,Q').
\]
Subtracting the layer representations and bounding the integral therefore yields
\[
\begin{aligned}
|(KQ)(E)-(KQ')(E)|
&=
\left|
\int_{(0,1]}
\left[
Q\{D:\eta_{\mathrm{layer}}\le\varphi_{K,E}(D)\}
-
Q'\{D:\eta_{\mathrm{layer}}\le\varphi_{K,E}(D)\}
\right]d\eta_{\mathrm{layer}}
\right|\\
&\le \int_{(0,1]}d_{\mathrm{TV}}(Q,Q')\,d\eta_{\mathrm{layer}}
=d_{\mathrm{TV}}(Q,Q').
\end{aligned}
\]
Taking the supremum over \(E\) proves
\[
d_{\mathrm{TV}}(KQ,KQ')\le d_{\mathrm{TV}}(Q,Q').
\]
% lean: Causalean.Stat.tvDist_integral_range
% lean: CausalSmith.Stat.PrivateCateRoughdesign.kernel_TV_contraction

\item Assume \(Q\ll Q'\) and the stated squared-deviation integral is finite. Choose an extended nonnegative Radon--Nikodym derivative and define, on all of \(\mathcal O^n\),
\[
p_{\mathrm{RN}}(D)=
\begin{cases}
\dfrac{dQ}{dQ'}(D),&\dfrac{dQ}{dQ'}(D)<\infty,\\
0,&\dfrac{dQ}{dQ'}(D)=\infty,
\end{cases}
\qquad
g_{\mathrm{RN}}(D)=p_{\mathrm{RN}}(D)-1.
\]
The infinite-density set is \(Q'\)-null, so this convention preserves the Radon--Nikodym integral identities. Finiteness of the nonnegative squared-deviation integral gives
\[
g_{\mathrm{RN}}\in L^2(Q'),
\qquad
\int g_{\mathrm{RN}}(D)^2\,Q'(dD)=\chi^2(Q,Q').
\]
Both \(|g_{\mathrm{RN}}|\) and the constant function one belong to \(L^2(Q')\). Hölder's inequality with exponents two and two gives
\[
\begin{aligned}
\int|g_{\mathrm{RN}}(D)|\,Q'(dD)
&=\int|g_{\mathrm{RN}}(D)|\cdot1\,Q'(dD)\\
&\le
\left(\int|g_{\mathrm{RN}}(D)|^2\,Q'(dD)\right)^{1/2}
\left(\int1^2\,Q'(dD)\right)^{1/2}\\
&=\sqrt{\chi^2(Q,Q')}.
\end{aligned}
\]
% lean: CausalSmith.Stat.PrivateCateRoughdesign.TV_le_sqrt_chiSq_of_finite
% lean: Causalean.Stat.tvDist_le_half_sqrt_chiSqDiv

To relate this integral to total variation, the probability normalization and the Radon--Nikodym identity give
\[
\int g_{\mathrm{RN}}\,dQ'
=
\int p_{\mathrm{RN}}\,dQ'-1
=Q(\mathcal O^n)-1=0.
\]
For every measurable dataset event \(F\),
\[
Q(F)-Q'(F)=\int_F g_{\mathrm{RN}}\,dQ',
\qquad
\int_{F^c}g_{\mathrm{RN}}\,dQ'
=-\int_F g_{\mathrm{RN}}\,dQ'.
\]
The pointwise inequalities
\(g_{\mathrm{RN}}\le|g_{\mathrm{RN}}|\) and
\(-g_{\mathrm{RN}}\le|g_{\mathrm{RN}}|\), integrated over \(F\) and \(F^c\), imply
\[
\left|\int_F g_{\mathrm{RN}}\,dQ'\right|
\le\int_F|g_{\mathrm{RN}}|\,dQ',
\qquad
\left|\int_F g_{\mathrm{RN}}\,dQ'\right|
=
\left|\int_{F^c}g_{\mathrm{RN}}\,dQ'\right|
\le\int_{F^c}|g_{\mathrm{RN}}|\,dQ'.
\]
Adding these bounds yields
\[
2|Q(F)-Q'(F)|
\le
\int_F|g_{\mathrm{RN}}|\,dQ'
+\int_{F^c}|g_{\mathrm{RN}}|\,dQ'
=
\int|g_{\mathrm{RN}}|\,dQ'.
\]
Taking the supremum over \(F\) and using the square-integrability bound proves
\[
d_{\mathrm{TV}}(Q,Q')
\le\frac12\int|g_{\mathrm{RN}}|\,dQ'
\le\frac12\sqrt{\chi^2(Q,Q')}.
\]
% lean: Causalean.Stat.measureReal_sub_eq_setIntegral_rnDeriv_sub_one
% lean: Causalean.Stat.abs_setIntegral_le_half_integral_abs_of_integral_eq_zero
% lean: Causalean.Stat.tvDist_le_half_integral_abs_rnDeriv
% lean: Causalean.Stat.tvDist_le_half_sqrt_chiSqDiv
% lean: CausalSmith.Stat.PrivateCateRoughdesign.private_kernel_comparison
\end{enumerate}
\end{proof}
